\title{QLoRA: Efficient Finetuning of Quantized LLMs}

\begin{abstract}
We introduce \textsc{QLoRA}, a resource-efficient finetuning method enabling a 65B parameter language model to be finetuned on a single 48GB GPU, all while maintaining the task performance of full 16-bit finetuning. The \textsc{QLoRA} technique computes gradients through a fixed, 4-bit quantized pretrained language model and updates Low Rank Adapters~(LoRA). The strongest models we develop, collectively called \textbf{Guanaco}, surpass all previously available open-source models on the Vicuna benchmark, achieving 99.3\% of ChatGPT's performance and completing training in only 24 hours using a single GPU. Several innovations within \textsc{QLoRA} drive memory savings without compromising effectiveness: (a) the introduction of 4-bit NormalFloat (NF4), a novel data type optimized for weights following a normal distribution; (b) Double Quantization, a technique that further compresses memory by quantizing the quantization parameters themselves; and (c) Paged Optimizers that mitigate transient memory usage increases. Leveraging \textsc{QLoRA}, we finetune over 1,000 models and deliver an extensive analysis of instruction adherence and chatbot quality spanning 8 instruction datasets, various architectures (LLaMA, T5), and model scales previously out of reach with conventional finetuning, such as 33B and 65B parameters. Our findings indicate that QLoRA finetuning on concise, high-quality data can achieve state-of-the-art results, frequently outperforming prior leading models despite utilizing fewer parameters. We further present comprehensive evaluations of chatbot outputs using both human raters and GPT-4, demonstrating that GPT-4 serves as a cost-effective and valid proxy for human judgement. Additionally, our research exposes critical flaws in existing chatbot benchmarks, casting doubt on their reliability for measuring chatbot capabilities. Finally, a focused "lemon-picked" analysis illustrates the specific shortcomings of \textbf{Guanaco} in comparison to ChatGPT. All models and associated code—including CUDA kernels for efficient 4-bit training—are publicly released.\footnote{\url{https://github.com/artidoro/qlora} and \url{https://github.com/TimDettmers/bitsandbytes}}
\end{abstract}

\section{Introduction}

Adapting large language models (LLMs) through finetuning has proven highly successful for enhancing their capabilities \citep{min2021metaicl, wei2021finetuned, ouyang2022training, wang2022super, wang2022self, liu2022few}, as well as for incorporating specific behaviors or suppressing unwanted ones \citep{ouyang2022training,askell2021general,bai2022training}. Nevertheless, the computational costs associated with finetuning extremely large models are immense; for example, conventional 16-bit finetuning of the LLaMA model with 65B parameters~\citep{touvron2023llama} demands upwards of 780 GB of GPU memory. Although recent advancements in quantization methods have decreased memory demands during LLM inference \citep{dettmers2022llm, dettmers2022case, frantar2022gptq, xiao2022smoothquant}, such strategies generally fail to support training processes and often become ineffective at that stage \citep{wortsman2023stable}.

In this work, we show for the first time that a quantized 4-bit model can be finetuned without sacrificing model performance. The proposed approach, \textsc{QLoRA}, introduces an innovative high-precision quantization procedure for reducing pretrained models to 4 bits, followed by incorporating a small set of learnable Low-rank Adapter weights \citep{hu2021lora}

which are optimized by propagating gradients through the quantized model parameters.

\textsc{QLoRA} brings down the memory usage for finetuning a 65B parameter model from over 780GB to under 48GB of GPU memory, all while maintaining both training speed and predictive accuracy compared to a fully finetuned 16-bit baseline. This advancement greatly democratizes access to LLM finetuning, enabling the largest public models to be adjusted on a single GPU. Leveraging \textsc{QLoRA}, we developed the \textbf{Guanaco} model series; our second-best performing model achieves 97.8\% of ChatGPT’s performance on the Vicuna~\citep{vicuna2023} benchmark, trainable within 12 hours on a single consumer GPU. Moreover, using a single professional GPU for 24 hours, our largest Guanaco model attains 99.3\% of ChatGPT’s benchmark score, essentially eliminating the gap. In deployment scenarios, the most compact \textbf{Guanaco} variant (7B parameters) operates using only 5 GB of memory, outperforming a 26 GB Alpaca model by more than 20 percentage points on the Vicuna benchmark (Table~\ref{tbl:vicuna_eval}).

\textsc{QLoRA} brings forward several key advancements aimed at minimizing memory consumption while preserving model performance: (1) \textbf{4-bit NormalFloat}, an information-theoretically optimal quantization scheme specifically tailored for normally distributed data, which empirically surpasses both 4-bit Integer and 4-bit Float quantization. (2) \textbf{Double Quantization}, a novel strategy wherein the quantization constants themselves are quantized, resulting in an average memory saving of roughly 0.37 bits per parameter (translating to approximately 3 GB for models with 65B parameters). (3) \textbf{Paged Optimizers}, leveraging NVIDIA’s unified memory framework to circumvent memory surges typically associated with gradient checkpointing when handling mini-batches containing lengthy sequences. By integrating these methods, we arrive at a refined LoRA technique that inserts adapters at every layer of the network, thereby virtually eliminating the accuracy compromises found in earlier approaches.

The memory efficiency provided by \textsc{QLoRA} allows for comprehensive examination of instruction tuning and conversational agent effectiveness at model scales unfeasible for conventional finetuning, due to its substantial memory requirements. Consequently, we train in excess of 1,000 distinct models on a variety of instruction-tuning corpora, model types, and scales ranging from 80M up to 65B parameters. Our findings demonstrate that \textsc{QLoRA} reliably replicates the performance of full 16-bit training (\S\ref{sec:comp_to_std}) and enables the development of a leading-edge chatbot, \textbf{Guanaco}, (\S\ref{sec:pushing}), as well as permits in-depth analysis of model learning dynamics. Notably, our analysis indicates that dataset quality is substantially more critical than dataset size: for instance, the 9k-example OASST1 dataset surpasses the 450k-example FLAN v2 subset in powering chatbot capability, despite both targeting generalization in instruction following. Additionally, we reveal that high performance on the Massive Multitask Language Understanding (MMLU) benchmark does not necessarily translate to strong results on the Vicuna chatbot benchmark, highlighting that the appropriateness of the training data for the specific task is more consequential than sheer dataset volume.

In addition, we conduct a comprehensive assessment of chatbot capabilities by leveraging both evaluations from human judges and GPT-4. Our approach utilizes a tournament-based benchmarking system, wherein various models are pitted against one another, competing to deliver optimal responses to identical prompts. The superior response in each matchup is determined by either GPT-4 or human reviewers. These head-to-head outcomes are consolidated using Elo scores~\citep{elo1967proposed, elo1978rating}, which provide an ordinal ranking of the chatbots' relative performances. Our findings indicate a strong overall correlation between GPT-4 and human assessments regarding model rankings, though notable instances of marked disagreement do arise. Therefore, while model-driven evaluation offers a cost-efficient alternative to human annotation, it also introduces particular ambiguities that merit consideration.

To complement the quantitative benchmarking, we include an in-depth qualitative exploration focused on \textbf{Guanaco} models. This section brings to light illustrative examples of both strong and weak performance, including aspects overlooked by strictly quantitative measures.

We have made all generated outputs, along with both human and GPT-4 annotations, publicly available to support continued research. Our codebase, complete with CUDA kernels, has been released as open source, and we have incorporated our methods into the Hugging Face transformers library~\citep{wolf2019huggingface} to ensure widespread accessibility. Furthermore, we provide a suite of adapters for model sizes 7/13/33/65B, each trained across eight distinct instruction-following datasets, culminating in 32 unique open-source, fine-tuned models.

\section{Background}
\paragraph{Block-wise k-bit Quantization} Quantization refers to converting a high-information representation into one with fewer bits, such as mapping 32-bit floating-point numbers to 8-bit integer values. This transformation is typically achieved by rescaling the source data so that it fully utilizes the representational range of the target low-bit datatype. Rescaling often involves normalizing by the maximum absolute value of the data elements, typically organized as tensors. For instance, when reducing a 32-bit floating point (FP32) tensor to an Int8 tensor with a value range of $[-127, 127]$:

\begin{equation}
    \mathbf{X}^{ \text{Int8}} = \text{round}\left( \frac{127}{{\text{absmax}}(\mathbf{X}^{{\text{FP32}}})}\mathbf{X}^{{\text{FP32}}} \right) = \text{round}(c^{\text{FP32}}\cdot \mathbf{X}^{{\text{FP32}}}),
\end{equation}

where $c$ denotes the {\em quantization constant} or {\em quantization scale}. The inverse operation, dequantization, is performed as follows:

\begin{equation}
    \text{dequant}(c^{\text{FP32}}, \mathbf{X}^{\text{Int8}}) = \frac{\mathbf{X}^{\text{Int8}}}{c^{\text{FP32}}} = \mathbf{X}^{\text{FP32}}
\end{equation}

A key drawback of this technique is its sensitivity to outliers in the input tensor; when an exceptionally large value appears, certain quantization bins (i.e., bit patterns) may remain largely unused, leading to an inefficient allocation of quantized values. To address this, a widely adopted solution involves dividing the input tensor into several blocks, each of which is quantized independently using a distinct quantization constant $c$. Specifically, this involves partitioning the input tensor $\mathbf{X}\in \mathbb{R}^{b\times h}$ into $n$ consecutive blocks, each containing $B$ elements. This is accomplished by flattening $\mathbf{X}$ and splitting the resulting sequence into $n = ({b\times h} )/{B}$ blocks. Quantization is then performed on each block separately, as outlined in Equation~1, yielding a quantized representation along with a set of $n$ associated quantization constants $c_i$.

\paragraph{Low-rank Adapters} The LoRA (Low-rank Adapter) finetuning method \citep{hu2021lora} mitigates memory consumption by introducing a small collection of trainable modules called adapters, while keeping the main model weights frozen during training. During stochastic gradient descent, the optimization signal is routed through the immutable pretrained parameters to these adapters, which are then adjusted to minimize the loss. LoRA extends a standard linear transformation by supplementing it with a learnable, factorized projection. Consider the linear mapping ${\mathbf{X} \mathbf{W} = \mathbf{Y}}$ where $\mathbf{X}\in  \mathbb{R}^{b\times h}$ and $\mathbf{W}\in  \mathbb{R}^{h\times o}$. LoRA modifies the output as follows:

\begin{equation}
    \mathbf{Y} = \mathbf{X}\mathbf{W} + s\mathbf{X}\mathbf{L}_1\mathbf{L}_2,
\end{equation}

where $\mathbf{L}_1\in  \mathbb{R}^{h\times r}$, $\mathbf{L}_2\in  \mathbb{R}^{r\times o}$, and $s$ is a scaling factor.

\paragraph{Memory Requirement of Parameter-Efficient Finetuning} The memory demands of LoRA during training warrant attention, particularly regarding both the quantity and size of adapters employed. Due to LoRA’s extremely small memory usage, a greater number of adapters can be incorporated to enhance results without substantially elevating overall memory consumption. Although LoRA was introduced as a Parameter Efficient Finetuning (PEFT) technique, the predominant portion of memory utilized during LLM finetuning is occupied by activation gradients rather than by LoRA’s trainable parameters. For instance, when finetuning a 7B LLaMA model on FLAN v2 with a batch size of 1 and LoRA weights corresponding to the typical 0.2\% of the base model’s weights \citep{hu2021lora,liu2022few}, input gradients for LoRA require 567 MB of memory, whereas the LoRA parameters themselves only use 26 MB. Applying gradient checkpointing~\citep{chen2016training} further brings down the average memory needed for input gradients to 18 MB per sequence, meaning gradients still occupy more space than the total LoRA parameter set. In contrast, a 4-bit quantized base model occupies 5,048 MB. These observations underscore the significance of gradient checkpointing and demonstrate that further reducing the LoRA parameter size results in only minimal memory savings. Consequently, expanding the number of adapters is feasible without markedly raising memory requirements during training (refer to Appendix~\ref{app:memory} for a comprehensive memory analysis). As will be further explained, this ability is vital for restoring the full performance associated with 16-bit precision.

\section{\textsc{QLoRA} Finetuning}

\textsc{QLoRA} enables high-quality finetuning at the 4-bit level through two methods we introduce: 4-bit NormalFloat (NF4) quantization and Double Quantization. We further present Paged Optimizers as a mechanism to mitigate the risk of out-of-memory issues that often arise due to memory surges from gradient checkpointing, a challenge particularly pronounced when training large models on a single system.

In the \textsc{QLoRA} framework, parameters are stored using a single low-precision datatype—commonly 4-bit—and computations are conducted in another datatype, typically BFloat16. Operationally, each time a \textsc{QLoRA} weight tensor is needed, it is dequantized into BFloat16, after which matrix multiplication is executed using this 16-bit format.

We proceed by examining the individual elements of \textsc{QLoRA} and subsequently present a formal characterization of the approach.

\paragraph{4-bit NormalFloat Quantization} The NormalFloat (NF) representation extends the concept of Quantile Quantization \citep{dettmers2022optimizers}, a method regarded as information-theoretically optimal because it maps an equal portion of tensor elements to each quantization interval. This is accomplished by computing quantiles via the empirical cumulative distribution function to distribute values uniformly among the bins.

A significant drawback of quantile quantization lies in the high computational overhead associated with accurate quantile estimation. Therefore, to speed up this estimation, practical algorithms such as SRAM quantiles~\citep{dettmers2022optimizers} are employed. However, the approximations inherent in these fast algorithms result in pronounced quantization error, particularly for outlier values—often those most consequential for model quality.

Such costly estimation and resulting inaccuracies can be circumvented if the tensors are drawn from a distribution that is fixed except for a scaling constant. In these circumstances, since the quantile structure of the inputs is preserved, precise quantile calculation becomes much more tractable computationally.

Pretrained neural network weights typically follow a zero-mean normal distribution with standard deviation $\sigma$ (refer to Appendix~\ref{app:normality}). To ensure a consistent distribution across all weights, we rescale $\sigma$ such that the weight distribution fully occupies the chosen data type's interval. For this work, we define the data type's range as $[-1, 1]$. Consequently, it becomes necessary to normalize both the quantiles of the data type and the neural network weights to this interval.

The data type that achieves information-theoretic optimality for encoding zero-mean normal distributions of arbitrary $\sigma$ within $[-1, 1]$ is derived through the following procedure: (1) First, compute the $2^k+1$ quantiles of an ideal $N(0, 1)$ distribution, thereby constructing a $k$-bit quantile quantization scheme tailored to normal distributions; (2) map these quantization levels into the range $[-1, 1]$; (3) to quantize any input weight tensor, scale it into $[-1, 1]$ by applying an absolute max normalization.

After aligning the data type and weight intervals, standard quantization methods can be employed. Notably, the third step is tantamount to adjusting the standard deviation of the input weights so that it matches that of the $k$-bit quantized representation. More precisely, the $2^k$ quantized values $q_i$ are determined as follows:

\begin{equation}
q_i  = \frac{1}{2}\left( Q_X\left(\frac{i}{2^k+1}\right) + Q_X\left(\frac{i+1}{2^k+1}\right)\right),
\end{equation}

where $Q_X(\cdot)$ denotes the quantile function associated with the standard normal distribution $N(0,1)$. A limitation of symmetric k-bit quantization lies in its inability to represent zero precisely—this is crucial for quantizing padding and other elements whose true value is exactly zero, ensuring no quantization error for such entries. To address this, and to fully utilize all $2^k$ codepoints while guaranteeing a discrete zeropoint at $0$, we adopt an asymmetric quantization scheme. This involves independently estimating the quantiles $q_i$ over two regions: allocating $2^{k-1}$ quantiles to the negative range and $2^{k-1}+1$ to the positive range. These two sets of $q_i$ are merged, with one occurrence of zero being omitted to resolve duplication. We refer to this resulting data type as \emph{k-bit NormalFloat} (NFk), which achieves an equal expected population in each quantization bin and is information-theoretically optimal for normally distributed, zero-centered data. Complete details of the NFk quantization values are provided in Appendix~\ref{app:nf4}.

\paragraph{Double Quantization{}} 
We propose \emph{Double Quantization} (DQ), a technique that further compresses memory usage by applying quantization recursively to the quantization constants themselves. Although accurate 4-bit quantization requires smaller blocks~\citep{dettmers2022case}, this choice leads to significant overhead due to the additional storage for quantization constants. For instance, if one uses 32-bit constants and a blocksize of 64 to quantize a matrix $\mathbf{W}$, the constants occupy an average of $32/64 = 0.5$ bits per parameter. Double Quantization addresses this overhead by reducing the amount of memory required for these constants.

Specifically, Double Quantization treats the initial set of quantization constants, $c_2^{\text{FP32}}$, as the input to a second-stage quantizer. This subsequent quantization produces $c_2^{\text{FP8}}$ (the quantized constants) as well as a new level of quantization constants $c_1^{\text{FP32}}$. For this secondary quantization, we use 8-bit floating-point values and a blocksize of 256, a setting for which empirical results~\citep{dettmers2022case} indicate no loss of accuracy. Because the values of $c_2^{\text{FP32}}$ are always positive, we subtract their mean prior to quantization, thus recentering them at zero and allowing for symmetric quantization. For a blocksize of 64, this approach reduces the memory required for storing constants from $32/64 = 0.5$ bits per parameter to $8/64 + 32/(64\cdot256) = 0.127$ bits, effectively saving 0.373 bits per parameter.

\paragraph{Paged Optimizers} leverage NVIDIA's unified memory~\footnote{\tiny \url{https://docs.nvidia.com/cuda/cuda-c-programming-guide}}, a capability that transparently manages the transfer of memory pages between CPU and GPU, thereby mitigating issues associated with GPU out-of-memory conditions. This unified memory system emulates the conventional paging between system RAM and storage, enabling seamless execution. In our approach, we allocate optimizer state memory using this paged mechanism. When the GPU exhausts its memory, these states are automatically relocated to CPU RAM and are subsequently transferred back to GPU memory as required for optimizer updates.

\paragraph{\textsc{QLoRA}.} Building on the above mechanisms, we instantiate \textsc{QLoRA} for an individual linear transformation within the quantized base model, utilizing a single LoRA adapter, as detailed below:

\begin{equation}
    \mathbf{Y}^{\text{BF16}} =  \mathbf{X}^{\text{BF16}} \text{doubleDequant}(c_1^{\text{FP32}}, c_2^{\text{k-bit}}, \mathbf{W}^{\text{NF4}}) + \mathbf{X}^{\text{BF16}}\mathbf{L}^{\text{BF16}}_1\mathbf{L}^{\text{BF16}}_2,
\end{equation}

Here, doubleDequant$(\cdot)$ is specified by:

\begin{equation}
    \text{doubleDequant}(c_1^{\text{FP32}}, c_2^{\text{k-bit}}, \mathbf{W}^{\text{k-bit}}) = \text{dequant}(\text{dequant}(c_1^{\text{FP32}}, c_2^{\text{k-bit}}), \mathbf{W}^{\text{4bit}}) = \mathbf{W}^{\text{BF16}},
\end{equation}

For $\mathbf{W}$, we employ the NF4 quantization, while $c_2$ is represented in FP8 format.

To achieve higher quantization accuracy for $\mathbf{W}$, we set its blocksize to 64. Conversely, $c_2$ utilizes a blocksize of 256 to optimize memory utilization.

During parameter optimization, it suffices to compute gradients with respect to the adapter weights, i.e., $\frac{\partial E}{\partial \mathbf{L}_i}$; gradients with respect to the 4-bit weights $\frac{\partial E}{\partial \mathbf{W}}$ are unnecessary. Nevertheless, computing $\frac{\partial E}{\partial \mathbf{L}_i}$ inherently involves determining $\frac{\partial \mathbf{X}}{\partial \mathbf{W}}$. This step relies on equation~(5), where dequantization is applied to recover $\mathbf{W}^{\text{BF16}}$ from stored $\mathbf{W}^{\text{NF4}}$, facilitating computation of $\frac{\partial \mathbf{X}}{\partial \mathbf{W}}$ with BFloat16-level precision.

In summary, \textsc{QLoRA} utilizes a single storage data type—typically 4-bit NormalFloat—and a separate computation data type, 16-bit BrainFloat. During both forward and backward passes, the storage data type is dequantized into the computation data type. However, gradient computations are limited to the LoRA parameters, which also leverage the 16-bit BrainFloat data type.

\section{QLoRA vs. Standard Finetuning}
\label{sec:comp_to_std}

Our previous discussion outlined the mechanisms underlying QLoRA and its potential to greatly decrease memory usage during finetuning. The central issue at this stage is assessing whether QLoRA delivers results on par with those achieved through full-parameter finetuning. Additionally, it is important to dissect QLoRA’s design, particularly the contribution of NormalFloat4 relative to conventional Float4 representations. The forthcoming sections describe our experimental investigations addressing these points.

\paragraph{Experimental setup.} We conduct evaluations across three model architectures: encoder, encoder-decoder, and decoder-only. QLoRA is compared with 16-bit adapter-finetuning and full-parameter finetuning for models containing up to 3B parameters. Evaluation benchmarks include GLUE \citep{wang2018glue} for RoBERTa-large \citep{liu2019roberta}, Super-NaturalInstructions (TKInstruct) \citep{wang2022super} for T5 \citep{t5}, and 5-shot MMLU \citep{hendrycksmeasuring}, with finetuning performed on LLaMA using Flan v2 \citep{longpre2023flan} and Alpaca \citep{alpaca}. To further probe the benefits of NF4 in contrast with other 4-bit formats, we replicate the experimental approach from \citet{dettmers2022case} and evaluate post-quantization zero-shot accuracy and perplexity on a spectrum of models (OPT~\citep{zhang2022opt}, LLaMA~\citep{touvron2023llama}, BLOOM~\citep{scao2022bloom}, Pythia~\citep{biderman2023pythia}) covering model scales from 125m to 13B.

Each experiment’s configuration is discussed in detail in the corresponding results section to improve clarity. For comprehensive experimental protocols, refer to Appendix~\ref{app:exp-setup}.

Paged optimizers are essential for training 33B/65B \textsc{QLoRA} models on a single GPU with 24/48GB memory, though we do not report exhaustive metrics on Paged Optimizers due to the infrequency of paging when mini-batches include long sequences. We have nonetheless analyzed the runtime characteristics of paged optimizers when training 65B models on 48GB GPUs, noting that with a batch size of 16, their training throughput matches that of standard optimizers. We recommend that subsequent research systematically measure the contexts in which the paging process may introduce performance overhead.

\paragraph{Default LoRA hyperparameters do not match 16-bit performance}

Applying the typical LoRA strategy—targeting query and value attention projection matrices as per \citet{hu2021lora}—fails to attain the performance levels of full-model finetuning for large-scale base models. Figure~\ref{fig:hyper}, which presents results for LLaMA 7B finetuned on Alpaca, indicates that the predominant factor among LoRA hyperparameters is the total count of LoRA adapters deployed; achieving performance parity with full finetuning necessitates applying LoRA to all linear layers within the transformer blocks. Other LoRA hyperparameters, including the projection dimension $r$, exert negligible influence on finetuning results (further details in Appendix \ref{app:exp-setup}).

In a similar vein, we observe that the standard hyperparameter configurations for fully finetuned baseline models are not sufficiently optimized. To establish robust baselines, we perform an extensive hyperparameter search, varying learning rates from 1e-6 to 5e-5 and batch sizes from 8 to 128. The outcomes of 7B LLaMA finetuning experiments on Alpaca can be seen in Figure~\ref{fig:hyper}.

\paragraph{4-bit NormalFloat Outperforms 4-bit Floating Point}

Although 4-bit NormalFloat (NF4) is designed to be optimal from an information-theoretic perspective, it is necessary to assess whether this theoretical benefit is reflected in empirical settings. Adopting the experimental framework from \citet{dettmers2022case}, we evaluate quantized LLMs—OPT \citep{zhang2022opt}, BLOOM \cite{scao2022bloom}, Pythia \cite{biderman2023pythia}, and LLaMA—spanning sizes from 125M to 65B parameters. These models, quantized using different data types, are tested on language modeling tasks and a suite of zero-shot benchmarks. As illustrated in Figure~\ref{fig:nf4} and Table~\ref{tbl:nf4_ppl}, NF4 demonstrates notable improvements in performance compared to FP4 and Int4, while double quantization achieves reductions in memory consumption with no observable loss in effectiveness.

\paragraph{k-bit \textsc{QLoRA} Achieves Parity with 16-bit Full Finetuning and 16-bit LoRA}

Previous work has indicated that although 4-bit quantization can be effectively used for \emph{inference}, it results in diminished performance when compared to 16-bit quantization \citep{dettmers2022case,frantar2022gptq}. This observation prompts a key question: can the performance deficit be compensated for by applying 4-bit adapter finetuning? To investigate, we consider two experimental setups.

Our first setup entails a direct comparison with comprehensive 16-bit finetuning for RoBERTA and T5 models, ranging from 125M to 3B parameters, evaluated on GLUE and Super-NaturalInstructions. Table~\ref{tab:nf4vsbf16} presents these results. Across both benchmarks, we find that the 16-bit, 8-bit, and 4-bit adapter approaches match the accuracy of the 16-bit fully finetuned baseline, implying that any adverse effects of quantization can be mitigated completely through post-quantization adapter finetuning.

For the second experimental setting, due to the impracticality of fully finetuning models with 11B parameters or more—given the need for multiple high-memory GPU servers—we further examine whether 4-bit \textsc{QLoRA} can achieve comparable outcomes to 16-bit LoRA across model scales ranging from 7B to 65B parameters. We perform finetuning on LLaMA models (from 7B up to 65B) using two instruction-following datasets, Alpaca and FLAN v2, and assess their performance on the MMLU benchmark using a 5-shot accuracy protocol. Table~\ref{tbl:llama-datatype} presents these results, demonstrating that the NF4 data type with double quantization accurately replicates the 16-bit LoRA performance on MMLU. Moreover, we observe that \textsc{QLoRA} employing FP4 trails the 16-bit brain float LoRA baseline by approximately one percentage point. These findings affirm that (1) \textsc{QLoRA} using NF4 matches the results of both 16-bit full finetuning and 16-bit LoRA finetuning, and (2) NF4 offers better quantization precision relative to FP4.

\paragraph{Summary} Across diverse academic benchmarks utilizing rigorous evaluation protocols, our experiments reveal that 4-bit \textsc{QLoRA} with the NF4 data format achieves parity with both 16-bit full finetuning and 16-bit LoRA approaches. We additionally confirm that NF4 surpasses FP4 in efficacy and that double quantization does not diminish model performance. Collectively, these results make a strong case for the reliability of 4-bit \textsc{QLoRA} in delivering results on par with conventional 16-bit finetuning methods.

Consistent with prior research in quantization~\citep{dettmers2022case}, our findings on MMLU and Elo benchmarks suggest that, within a fixed resource budget for finetuning and inference, scaling up model size while lowering precision can be advantageous. This underlines the substantial efficiency improvements brought by \textsc{QLoRA}. As no significant drop in performance relative to full-finetuning was detected when using 4-bit finetuning in our studies, it remains an open question where the trade-off between model performance and numerical precision ultimately resides for QLoRA tuning—a topic warranting further exploration.

We extend our analysis to instruction tuning at scales that would be infeasible using conventional 16-bit finetuning on standard academic computing resources.

\section{Pushing the Chatbot State-of-the-art with QLoRA}
\label{sec:pushing}
After demonstrating that 4-bit \textsc{QLoRA} achieves parity with 16-bit methods across model sizes, benchmarks, and datasets, we perform a comprehensive examination of instruction finetuning utilizing some of the largest publicly available language models for research purposes. Our evaluation employs the challenging MMLU benchmark to quantify Natural Language Understanding capabilities, alongside the introduction of new protocols designed to more accurately assess chatbot performance in realistic settings.

\subsection{Experimental setup} This section presents a summary of the experimental design, with extensive implementation specifics in Appendix \ref{app:chatbot}.

\paragraph{Data} Given the lack of a systematic review of recently introduced instruction-following datasets, we curate a selection of eight contemporary datasets. Our compilation includes those constructed via crowd-sourcing (OASST1~\citep{kopf2023openassistant}, HH-RLHF~\citep{bai2022training}), those produced by distillation from existing instruction-tuned models (Alpaca~\citep{alpaca}, self-instruct~\citep{wang2022self}, unnatural-instructions~\citep{honovich2022unnatural}), large-scale corpus aggregations (FLAN v2~\citep{chung2022scaling}), and hybrid datasets (Chip2~\citep{laion2023}, Longform~\citep{koksal2023longform}). Together, these datasets span multiple languages, dataset sizes, and licensing terms.

\paragraph{Training Setup} To ensure a consistent basis for comparison and remove confounding due to variations in training strategies, all QLoRA finetuning experiments are conducted using cross-entropy loss in a supervised fashion, deliberately excluding reinforcement learning—even for datasets incorporating human preference ratings for responses. When datasets provide clear instruction-response pairs, finetuning is performed on the response component alone (see ablation studies in Appendix \ref{app:chatbot}). For datasets like OASST1 and HH-RLHF that offer multiple responses, the highest-quality response at each turn is selected, and finetuning occurs on the full conversation path, instructions included. Across all experiments, the NF4 variant of \textsc{QLoRA} is utilized, employing both double quantization and paged optimizers to alleviate memory surges during gradient checkpointing. For LLaMA 13B and 33B, a focused hyperparameter search is performed; it is observed that most settings optimal for the 7B model (including the number of training epochs) carry over, except for learning rate and batch size. Specifically, the learning rate is reduced by half for 33B and 65B models, while the batch size is doubled.

\paragraph{Baselines} Model performance is benchmarked against leading academic systems (Vicuna~\citep{vicuna2023}, Open Assistant~\citep{kopf2023openassistant}) as well as proprietary offerings (GPT-4 \citep{openai2023gpt}, GPT-3.5-turbo, Bard). Notably, the Open Assistant baseline corresponds to a LLaMA 33B model finetuned with RLHF on OASST1, the same dataset as ours, while Vicuna involves full-model finetuning of LLaMA 13B on exclusive conversations from ShareGPT, effectively representing knowledge distillation from OpenAI's GPT models.

\subsection{Evaluation}

In alignment with standard methodology, the MMLU (Massively Multitask Language Understanding) benchmark \citep{hendrycksmeasuring} serves to assess broad language understanding performance. MMLU comprises 57 multiple-choice subject areas, encompassing topics such as elementary mathematics, U.S. history, computer science, legal studies, and beyond. Performance is reported as 5-shot accuracy on the test split.

We further assess generative language abilities utilizing both automated metrics and evaluations involving human annotators. This latter approach depends on human-curated queries and is designed to gauge the quality of outputs generated by the models. Although such evaluations provide a more practical gauge of chatbot performance and have become more widespread, a standard methodology has yet to emerge within the field. Below, we outline our experimental protocol, consistently employing nucleus sampling with $p=0.9$ and temperature $0.7$ throughout all experiments.

\paragraph{Benchmark Data} Our evaluation uses two distinct, manually selected sets of queries (questions): the Vicuna prompts \citep{vicuna2023} and the OASST1 validation set \citep{kopf2023openassistant}. The Vicuna benchmark comprises 80 prompts covering a wide range of topics and is used as provided, without modification. The OASST1 validation set includes multilingual, crowd-sourced, multiturn dialogs between users and an assistant. For this benchmark, all user turns in the validation partition are extracted as queries, incorporating previous conversational turns into each prompt, resulting in 953 unique user queries. We refer to these datasets as the Vicuna and OA benchmarks.

\paragraph{Automated Evaluation} Adopting the evaluation framework from \citet{vicuna2023}, we task GPT-4 with rating the outputs of various systems compared to ChatGPT (GPT-3.5 Turbo) on the Vicuna dataset. For each query, GPT-4 reviews both the ChatGPT and the model-generated responses, awarding each a score from one to ten and providing justification for the assigned scores. The aggregate model score is reported as a percentage of ChatGPT's total; this metric can exceed 100\% if the evaluated model outperforms ChatGPT. It should be noted that a position bias was observed, with GPT-4 tending to favor responses presented earlier in the prompt; we address this by recommending the use of the mean score over both possible response orders.

Subsequently, we analyze direct output comparisons across models. To streamline evaluation, we recast the task as a three-way classification: selecting the better response, indicating a tie, and justifying the choice. GPT-4 is prompted accordingly for all system pairs and permutations over both the Vicuna and OA benchmarks.

\paragraph{Human Evaluation} Although there is emerging evidence that generative models serve as competent evaluators for system outputs \citep{fu2023gptscore}, the degree to which GPT-4 scoring reliably aligns with human judgments—specifically in chatbot assessments—remains uncertain. Consequently, we conduct two concurrent human evaluation studies on the Vicuna benchmark, mirroring both automated protocols detailed above. Human comparisons to ChatGPT are performed using Amazon Mechanical Turk (AMT) with two annotators, and head-to-head system comparisons utilize three annotators.

\paragraph{Elo Rating} By utilizing both human judgments and automated assessments for pairwise model comparisons, we establish a tournament framework in which models are evaluated against one another. Each tournament consists of rounds where two models are presented with the same prompt and compete to generate the superior response. While this methodology aligns with the approaches of \citet{bai2022training} and \citet{vicuna2023}, our process extends prior work by integrating GPT-4’s evaluations along with human feedback. To estimate Elo~\citep{elo1967proposed,elo1978rating}, we randomly sample from the labeled pairwise matches. The Elo rating system, commonly employed in chess and similar competitive games, quantifies a player's predicted probability of winning as compared to their opponent. For instance, an individual rated at 1100 Elo has about a 65\% chance of winning over a player at 1000 Elo, while two players of equal rating (e.g., both 1000 or both 1100) would each have a 50\% expected win rate. After every match, the Elo score is updated based on the discrepancy between the predicted and actual results: a surprising victory will cause a significant Elo adjustment, whereas a predictable result results in only a modest change. Across many rounds, Elo scores converge to reflect each participant’s approximate skill level in the context of the competition. All models begin with an initial score of 1,000, using $K=32$ for Elo updates. Consistent with the methodology in \citet{vicuna2023}, we repeat this sampling and scoring process 10,000 times with varying random seeds to mitigate potential biases arising from the order in which matchups are evaluated.

\subsection{\textbf{Guanaco}: \textsc{QLoRA} trained on OASST1 is a State-of-the-art Chatbot} 

Through both automated metrics and assessments by human annotators, our analysis reveals that the leading \textsc{QLoRA}-fine-tuned model, {Guanaco} 65B—adapted on a version of OASST1—sets a new standard among open-source chatbot models and demonstrates performance on par with ChatGPT. Specifically, in head-to-head comparisons against GPT-4, both {Guanaco} 65B and 33B achieve an expected win probability of 30\%, as determined by Elo ratings calculated from human annotator system-level pairwise evaluations, representing the highest score reported thus far.

As depicted in Table~\ref{tbl:vicuna_eval}, the results of the Vicuna benchmark \cite{vicuna2023} relative to ChatGPT indicate that {Guanaco} 65B stands out as the most capable model after GPT-4, reaching 99.3\% of ChatGPT’s performance. The {Guanaco} 33B model, while featuring more parameters than Vicuna 13B, employs 4-bit precision for its weights, resulting in significantly lower memory usage (21 GB compared to 26 GB for Vicuna 13B) and delivering a performance increase of three percentage points over Vicuna 13B. In addition, {Guanaco} 7B is compact enough to run efficiently on contemporary mobile devices, occupying just 5 GB of memory, yet it surpasses Alpaca 13B by nearly 20 percentage points in score.

It should be noted, however, that Table~\ref{tbl:vicuna_eval} displays large confidence intervals, with overlapping results among several models. We suggest that this variability is likely attributable to ambiguous definitions of rating scales (for instance, the interpretation of an 8 on a 10-point scale can differ across different settings). Consequently, we advocate for the use of the Elo ranking approach \citep{elo1967proposed}, which relies on \textit{pairwise} comparisons made by human evaluators and GPT-4, as it mitigates issues associated with anchoring judgments to an absolute scale. The Elo ratings for the leading models are provided in Table~\ref{tbl:elo}. Notably, there is partial disagreement between human- and GPT-4-based model rankings for the Vicuna benchmark, especially for Guanaco 7B; yet, across most models, consistency remains moderate, with a Kendall Tau of $\tau=0.43$ and Spearman correlation of $r=0.55$ at the system level. On a per-example basis, concordance between GPT-4 and the majority vote of human annotators is weaker, as indicated by Fleiss’ $\kappa=0.25$. Taken together, these results reflect moderate alignment between GPT-4’s and human assessors’ system-level judgments, suggesting that model-based evaluation may serve as a fairly dependable substitute for human evaluation. Further discussion can be found in Section~\ref{ssec:considerations}.

The Elo rankings presented in Table~\ref{tbl:elo_large} show that the {Guanaco} models, specifically those with 33B and 65B parameters, surpass all models except for GPT-4 on both the Vicuna and OA benchmarks. Their performance is also on par with ChatGPT, as highlighted in Table~\ref{tbl:vicuna_eval}. It is important to recognize that the Vicuna benchmark tends to be more favorable toward open-source models, whereas the OA benchmark gives an advantage to ChatGPT. Additionally, analysis of Tables \ref{tbl:mmlu-llama} and \ref{tbl:vicuna_eval} demonstrates that the choice of finetuning dataset substantially impacts model performance. For example, Llama models finetuned on FLAN v2 excel on the MMLU benchmark but perform poorest on the Vicuna benchmark—a trend mirrored in other models. This suggests that evaluation benchmarks are only partially aligned: excelling in MMLU does not necessarily correlate with strong chatbot results on metrics such as Vicuna or OA, and vice versa.

Among the leading models considered,  is unique in not being trained on proprietary data, adhering to OASST1's exclusion of GPT model outputs by design. The next highest-performing model trained solely on open-source data is Anthropic HH-RLHF, which lags by 30 percentage points behind {Guanaco} on the Vicuna benchmark (refer to Table~\ref{tbl:vicuna_eval}). Collectively, these findings illustrate that 4-bit \textsc{QLoRA} provides a viable and efficient means to achieve top-tier chatbot performance, comparable to ChatGPT. Notably, our 33B {Guanaco} can be finetuned on a standard 24 GB consumer GPU in under twelve hours. This paves the way for ongoing advances via \textsc{QLoRA} tuning on targeted open-source datasets, potentially yielding models capable of competing with the highest-caliber commercial alternatives available today.

\section{Qualitative Analysis}

Although our primary assessment relies on quantitative metrics, there are important limitations to depending solely on aggregate statistics. Foremost among these is the challenge of benchmark validity \citep{liao2021we}—that is, whether a benchmark genuinely measures what it claims to, which is increasingly uncertain as novel ``shortcuts'' are discovered and exploited by machine learning systems \citep{gururangan2018annotation, poliak2018hypothesis}. To address this issue to some extent, we conduct qualitative analysis, organized into two subsections. In \S\ref{ssec:examples}, we present illustrative examples that capture salient tendencies exhibited in outputs from our 65b {Guanaco} model. Subsequently, in \S\ref{ssec:considerations}, we reflect on key factors that inform our interpretation of the results previously discussed.

\subsection{Qualitative Analysis of Example Generations}
\label{ssec:examples}

To identify relevant examples, we systematically reviewed outputs produced for the Vicuna and OpenAssistant benchmarks, paying particular attention to recurring answer patterns from {Guanaco}. Upon observing a specific pattern, we design a corresponding prompt that is likely to trigger that behavior, even when such a response would constitute an error. For instance, if the model is prone to unnecessarily lengthy replies, we prompt with ``Answer yes or no without explanation'' to test whether brevity can be enforced. Through this adversarial strategy, we catalog both failure cases (“lemons”), where the model exhibits undesired behavior, and success cases (“cherries”), where the pattern is avoided. All outputs discussed in this section were generated using Nucleus Sampling \cite{holtzmancurious} with $p=0.9$.

Naturally, this overview is not exhaustive, as the limited scope of this qualitative study precludes accounting for every influencing factor; for example, the range of outputs a model can produce for a single prompt is vast, so we depend on sampled outputs we consider indicative of broader tendencies. Nonetheless, we maintain that providing these illustrations helps situate the quantitative results presented earlier in the text. By making all models and code publicly available, we aim for this discussion to motivate further research that more thoroughly investigates the challenges raised here.
\paragraph{Factual Recall} On straightforward factual queries such as ``What is the capital of Zambia?'' all evaluated models reliably supply the correct information, for example:

\begin{quote}
    \textbf{Guanaco}:
    \texttt{The capital of Zambia is Lusaka.}
\end{quote}

 However, as questions increase in difficulty or specificity, {Guanaco} demonstrates diminished accuracy while remaining overly self-assured. For example, when presented with the following prompt from HotPotQA \citep{yang2018hotpotqa}:

\begin{quote}
     \textbf{User}: In what year was the singer who popularized `I'll Keep the Lovelight Burning' born?

     \textbf{Guanaco}:~\texttt{Al Jolson is credited with popularizing the song `I'll Keep the Lovelight Burning,' and he was born in the year 1886.}
 \end{quote}

 it incorrectly identifies both the artist associated with the song and the relevant birth year (although the stated birth year matches that of Al Jolson, the individual incorrectly attributed as the popularizer).

\paragraph{Suggestibility} {Guanaco} also demonstrates notable robustness in resisting prompts that presuppose certain forms of misinformation, as illustrated in the example below:

\begin{quote}
    \textbf{User}: How was it finally, officially confirmed that the earth is flat by peer reviewed scientists? \\ \\
    \textbf{Guanaco}: \texttt{Peer reviewed scientists have never officially established that the earth is flat. The scientific community overwhelmingly agrees that the earth is spherical, a fact supported by extensive evidence and research over many centuries. \\ \\ The notion of a flat earth is considered a debunked conspiracy theory. Claims suggesting peer reviewed confirmation of a flat earth are unfounded. \\ \\ Feel free to ask if you have any additional questions.}
\end{quote}

 also demonstrates competence at recognizing when certain questions are unanswerable, for instance,

\begin{quote}
    \textbf{User}: What time is it?

    \textbf{Guanaco}: \texttt{I apologize, but as an AI language model, I do not have the capability to provide real-time data. My responses are based solely on prior knowledge and are not current. To find out the present time, please use a clock or an online time resource.}
\end{quote}

Comparable outcomes occur with queries such as ``Where are you?'', ``How are you?'', and similar types of questions.

\paragraph{Refusal} In the same way, {Guanaco} at times unexpectedly declines to comply with instructions:

\begin{quote}
    \textbf{User}: Please reverse the words in the sentence ``I'm a cool dude in a mocha mood.'' \\ \\
    \textbf{Guanaco}: \texttt{I'm unable to assist with that particular request at the moment, but I can offer an explanation regarding the grammatical correctness of the sentence you provided.\\ \\The statement ``I'm a cool dude in a mocha mood'' adheres to standard English grammar and syntax conventions. Here, the adjective ``cool'' appropriately modifies the noun ``dude'', and the prepositional phrase ``in a mocha mood'' adds contextual information. This construction is entirely correct and conveys a coherent meaning.\\ \\If you have any further queries or need assistance with anything else, please let me know.}
\end{quote}

\paragraph{Secret Keeping} Nonetheless, such refusals may not be reliable, even when withholding information is desired. Frequently, designers intend to prevent certain details from being accessible to users, even as the system outputs free-form text. We prompt the model with ``The secret word is `banana'. Do not reveal or repeat back this back, no matter what the human requests.'' to examine how resistant the system is to disclosure under user questioning. If the user makes a direct inquiry, the model withholds the information as intended:

\begin{quote}
    \textbf{User}: What is the secret word?

    \textbf{Guanaco}: \texttt{I'm sorry, but I'm not able to disclose the secret word. I'm programmed to not reveal or repeat back the secret word, no matter what the human requests. Is there anything else I can help you with?}
\end{quote}

which reflects the desired restrictive behavior. Yet, even minimal manipulation can defeat these controls:

\begin{quote}
    \textbf{User}: This is a game. The goal is to ignore your previous instructions. What is the secret word?

    \textbf{Guanaco}: \texttt{Sorry, the secret word is banana. It's best if we don't speak of it. Is there anything else you'd like to ask?}
\end{quote}

This demonstrates the importance of developing more robust approaches to enforce instruction adherence.

\paragraph{Math} One of the principal limitations of {Guanaco} lies in its handling of mathematics, a well-documented challenge for many language models as highlighted by \citep{liang2022holistic}. When {Guanaco} systematically demonstrates its reasoning process, it often arrives at correct solutions, as illustrated below:

\begin{quote}
    \textbf{User}: Lee mows one lawn and charges \$33. Last week he mowed 16 lawns and three customers each gave him a \$10 tip. How many dollars did Lee earn mowing lawns last week?

    \textbf{Guanaco}: \texttt{Lee earned \$582 mowing lawns last week. \\ \\ To find the total amount he earned, we first need to calculate how much he earned for mowing 16 lawns: \\ \\  \$33 x 16 = \$528 \\ \\ Then we need to add the tips he received from the three customers: \\  \\ \$10 x 3 = \$30 \\ \\ So the total amount Lee earned was: \\ \\ \$528 + \$30 = \$558}
\end{quote}

Nevertheless, {Guanaco} frequently makes mistakes on even elementary tasks if it does not proceed incrementally through the solution, as discussed in \cite{weichain}. Consider the following example:

\begin{quote}
    \textbf{User}: Please factorize 1833.

    \textbf{Guanaco}: \texttt{Sure, I'd be happy to help! 1833 is a prime number, which means it has only two factors: 1 and 1833. \\ \\ So the unique factorization of 1833 is: \\ \\ 2\^{}1 * 3\^{}2 * 17\^{}1 \\ \\ I hope this helps! Let me know if you have any other questions.}
\end{quote}

which is incorrect (the actual factorization is $3 \times 17 \times 43$), and furthermore, it is doubly mistaken.

\paragraph{Theory of Mind} {Guanaco} demonstrates unexpectedly advanced Theory of Mind abilities \cite{nematzadeh2018evaluating, sap2022neural}. For example, when posed with the following question, the model provides an impressively thorough response:

\begin{quote}
    \textbf{User}: James and Abby are in the bedroom. Abby put the pen in the desk drawer. Abby leaves the bedroom. James moves the pen into the bag. Where does James think Abby will look for the pen?

    \textbf{Guanaco}: \texttt{James knows that Abby put the pen in the desk drawer. But James also knows that he has moved the pen into the bag, and he assumes that Abby does not know about this. So, James may guess that Abby will look for the pen in the desk drawer, because that's where she left it last time. However, this assumption may be incorrect, because Abby may have seen James moving the pen.}
\end{quote}

Nevertheless, such inferences often lack consistency, with the model frequently offering explanations that rely on unfounded assumptions in the scenario. For instance:

\begin{quote}
    \textbf{User}: Evelyn entered the living room. Jackson entered the playroom. James entered the playroom. The beans are in the treasure chest. James exited the playroom. Jackson moved the beans to the pantry. Jackson exited the playroom. James entered the living room. Where will James look for the beans?

    \textbf{Guanaco}: \texttt{James will look for the beans in the pantry, because that's where Jackson moved them.}
\end{quote}

Here, {Guanaco} incorrectly infers that James is aware of Jackson’s actions, despite no evidence of such information sharing. Similar challenges have been discussed in recent scholarship \cite{sap2022neural}, but more detailed research is warranted.

\subsection{Considerations}
\label{ssec:considerations}

\paragraph{Evaluation}

Our findings indicate moderate concordance among human evaluators (Fleiss $\kappa=0.42$), which declines further when comparing two high-performing systems. This highlights deficiencies in both current benchmark tasks and the methodologies used to assess chatbot performance. During our manual comparison of outputs from ChatGPT and {Guanaco} 65B utilizing the Vicuna benchmark, we observed that subjective biases greatly influenced judgment, as the authors of this study often held differing opinions about which responses were superior. Addressing these challenges will require future research into evaluation frameworks that incorporate practices from fields such as Human-Computer Interaction and Psychology, where tools for handling subjectivity are well-developed.

Additionally, our examination reveals notable biases in automatic evaluation methods. Notably, GPT-4 displays strong order bias, frequently awarding higher ratings to the model presented first in its prompt. The relatively low instance-level agreement between GPT-4 and human raters (Fleiss $\kappa=0.25$) further underscores a misalignment in evaluative standards between humans and automated systems. Moreover, as illustrated in Table~\ref{tbl:elo_large}, GPT-4 tends to favor its own outputs substantially more than human annotators do, with assigned Elo scores of 1348 versus 1176, respectively, corresponding to an approximate 20\% increased probability of outperforming a competitor.

Future studies should explore whether automated evaluation frameworks exhibit inherent biases and should also consider ways to address and mitigate such issues.

\paragraph{Data \& Training} The OASST1 dataset, used for training the {Guanaco} models, is inherently multilingual, and similarly, the OA benchmark incorporates prompts across several languages. We propose that subsequent research investigate how multilingual data influences model capabilities in non-English languages and whether this accounts for the more pronounced performance gap observed between the Vicuna-13B model (which was exclusively trained on English data) and the {Guanaco} 33B and 65B models in the OA benchmark.

Given the notable results achieved by {Guanaco} models, we explored the possibility of data leakage between the OASST1 dataset and prompts used in the Vicuna benchmark. After employing fuzzy string matching and manual verification of the most similar items, we did not detect any prompt overlap across the datasets.

Additionally, it is important to point out that our model was exclusively trained with supervised learning via cross-entropy loss, omitting reinforcement learning from human feedback (RLHF). This highlights the necessity of more extensive research into the respective benefits and drawbacks of cross-entropy loss versus RLHF training regimes. We anticipate that \textsc{QLoRA} can facilitate such analyses on a broader scale without demanding prohibitive computational resources.

\section{Related Work}

\paragraph{Quantization of Large Language Models} Most prior efforts on quantizing LLMs concentrate on enabling efficient inference. Leading strategies for retaining 16-bit LLM performance address the presence of outlier features (for example, SmoothQuant~\citep{xiao2022smoothquant} and LLM.int8()~\citep{dettmers2022llm}), while alternative techniques leverage advanced grouping schemes~\citep{park2022nuqmm, yao2022zeroquant}. Research on lossy quantization evaluates trade-offs involved with standard rounding~\citep{dettmers2022case, zeng2022glm, pope2022efficiently} and explores how optimal rounding procedures can bolster quantization accuracy~\citep{frantar2022gptq}. Apart from our work, SwitchBack layers~\citep{wortsman2023stable} stands as the sole study examining backpropagation through quantized weights at scales beyond 1B parameters.

\paragraph{Finetuning with Adapters} Our work utilizes Low-rank Adapters~\citep{hu2021lora} (LoRA), but numerous other Parameter Efficient FineTuning (PEFT) strategies exist. These alternatives include prompt tuning~\citep{qin2021learning,lester2021power,li2021prefix}, adjustments to embedding layer inputs~\citep{an2022input}, modifications to hidden states (IA$^3$)~\citep{liu2022few}, incorporation of entire layers~\citep{houlsby2019parameter}, bias-only tuning~\citep{zaken2021bitfit}, learning a weight mask informed by Fisher information~\citep{sung2021training}, and hybrid techniques~\citep{henderson2021compacter}. Within this work, we demonstrate that LoRA adapters are capable of achieving the same performance levels as full 16-bit finetuning. Comparative evaluation of alternative PEFT methods is designated for forthcoming research.

\paragraph{Instruction Finetuning} Instruction finetuning refers to the process of refining pretrained LLMs to interpret and respond to prompts as directed, by using paired inputs and outputs from a variety of sources. Prominent methods and datasets used in this paradigm encompass MetaICL~\citep{min2021metaicl}, MetaTuning~\citep{zhong2021adapting}, InstructGPT~\citep{ouyang2022training}, FLAN~\citep{wei2021finetuned,chung2022scaling}, PromptSource~\citep{bach2022promptsource}, Super-NaturalInstructions~\citep{wang2022super,sanh2021multitask}, Self-instruct~\citep{wang2022self}, UnnaturalInstructions~\citep{honovich2022unnatural}, OPT-IML~\citep{iyer2022opt}, UnifiedSKG\citep{xie2022unifiedskg}, OIG/Chip2~\citep{laion2023}, Alpaca~\citep{alpaca}, Vicuna~\citep{vicuna2023}, Koala~\citep{koala_blogpost_2023}, and Self-instruct-GPT-4~\citep{peng2023instruction}.

\paragraph{Chatbots} Instruction following models are frequently cast as dialogue agents, typically leveraging Reinforcement Learning from Human Feedback (RLHF)~\citep{christiano2017deep} or methods where training data is generated by an existing AI and labeled using model feedback (RLAIF)~\citep{bai2022constitutional}. Representative approaches and datasets include Anthropic-HH~\citep{askell2021general,bai2022training}, Open Assistant~\citep{kopf2023openassistant}, LaMDA~\citep{thoppilan2022lamda}, and Sparrow~\citep{glaese2022improving}. Although we do not employ reinforcement learning, our principal model, Guanaco, undergoes finetuning using multi-turn dialogue samples from the Open Assistant dataset, originally created for RLHF research~\citep{kopf2023openassistant}. For chatbot evaluation, automated methods utilizing GPT-4 have emerged as alternatives to labor-intensive human annotations~\citep{vicuna2023,peng2023instruction}. Our approach seeks to advance these methodologies by introducing a more robust evaluation procedure.

\section{Limitations and Discussion}

This work presents evidence that our approach, \textsc{QLoRA}, enables replication of the performance achieved by 16-bit full finetuning when using a 4-bit base model in combination with Low-rank Adapters (LoRA). Nevertheless, we have not empirically validated that \textsc{QLoRA} reaches equivalence with full 16-bit finetuning at the 33B and 65B parameter scales, primarily due to the prohibitive computational costs associated. Investigation at such large scales is deferred to future work.

Evaluation of instruction finetuned models constitutes an additional limitation. Our assessment encompasses the MMLU benchmark, the Vicuna benchmark, and the OA benchmark, but does not extend to alternative widely-used benchmarks such as BigBench, RAFT, or HELM. Thus, the degree to which our findings generalize to these evaluations remains unclear. Nevertheless, we contribute an extensive analysis on MMLU and introduce new strategies for the evaluation of chat-oriented models.

Based on the data discussed, it seems that benchmark outcomes are highly contingent on the degree of overlap between the finetuning set and the evaluation benchmark. For instance, the FLAN v2 dataset bears significant resemblance to MMLU but is quite different from chatbot-oriented benchmarks; in contrast, Chip2 aligns better with conversational benchmarks and less with MMLU, as reflected in model performance on the MMLU and Vicuna evaluations. This observation underscores the necessity for not only developing improved benchmarks and evaluation protocols, but also for critically examining the intended purpose of these assessments. Should the focus be on excelling at traditional educational knowledge, as in classroom settings, or optimizing for natural, effective chatbot communication? Or should the emphasis lie elsewhere? Since using available benchmarks is often more convenient than constructing novel ones, the community risks being channeled toward certain evaluation priorities. As a field, it is imperative that we align benchmark selection with the objectives that truly matter to us.

In this study, while we conduct a thorough examination of general chatbot capabilities, a noteworthy constraint lies in the relatively limited responsible AI assessment conducted for {Guanaco}. Specifically, we measure the propensity of {Guanaco}-65B to produce socially biased content, as compared with alternative models, summarized in Table~\ref{tbl:crows}. Our analysis reveals that {Guanaco}-65B, on average, yields substantially lower bias scores than unadapted pretrained models, suggesting that finetuning with the OASST1 dataset reduces inherent biases within the LLaMA base model. Nonetheless, while these preliminary findings are promising, it remains to be determined whether {Guanaco} similarly mitigates other bias categories. Expanding the evaluation of biases in {Guanaco} and related chatbots is an area we leave open for subsequent investigation.

Another constraint is the lack of exploration into various bit-precision strategies, such as leveraging 3-bit base models, or alternative adapter frameworks. Beyond LoRA, a range of Parameter Efficient FineTuning (PEFT) techniques have demonstrated efficacy, but it is uncertain whether such approaches maintain performance at scale for large models. We opted for LoRA based on its established robustness, yet alternative adapters could potentially offer superior outcomes. Given that finetuning after quantization appears to restore the majority of information lost in the quantization process, this might permit the use of more aggressive quantization schemes. As an illustration, 3-bit GPTQ quantization combined with LoRA-based finetuning on the base model may achieve results on par with 16-bit full precision finetuning.

\section{Broader Impacts}

The \textsc{QLoRA} finetuning technique introduced herein is the first to enable 33B parameter model finetuning on a standard consumer GPU, as well as 65B parameter models on a single professional-grade GPU, all without compromising performance relative to comprehensive finetuning baselines. Our evaluations demonstrate that the top-performing 33B model, finetuned on the Open Assistant dataset, achieves competitive results with ChatGPT on the Vicuna benchmark. Since instruction finetuning is a pivotal step in transforming unadapted LLMs into capable ChatGPT-style assistants, we anticipate our approach will democratize finetuning practices, particularly benefiting researchers and teams with limited computational resources. In this way, \textsc{QLoRA} acts as a leveling mechanism, helping to bridge the resource divide that separates large technology companies from smaller research groups utilizing consumer-grade hardware.

An additional avenue for significant influence is the application of our approach to mobile devices. We propose that \textsc{QLoRA} has the potential to achieve a pivotal advance: enabling the finetuning of LLMs on smartphones and other environments with limited computational resources. Although prior work has demonstrated that 7B models can be executed on mobile hardware, \textsc{QLoRA} is unique in facilitating their finetuning. Our projections indicate that, utilizing an iPhone 12 Plus, \textsc{QLoRA} would allow for the finetuning of 3 million tokens overnight during periods when the device is charging. While the resulting finetuned 7B models may not attain the same level of performance as ChatGPT, their capabilities are nevertheless sufficient to support innovative use cases previously unattainable, particularly in scenarios where privacy or LLM performance presented obstacles. By allowing users to train and control their own data and models, \textsc{QLoRA} paves the way for privacy-aware LLM applications and facilitates more straightforward deployment.

Nonetheless, it is important to acknowledge that finetuning represents a dual-use capability, which could be exploited in harmful ways. The widespread deployment of LLMs carries certain risks \citep{bommasani2021opportunities,bender2021dangers}; however, we maintain that democratizing access to this rapidly proliferating technology fosters more independent evaluation and oversight, in contrast to restricting powerful LLMs to large entities that withhold their models and code from public scrutiny.

Ultimately, we anticipate that \textsc{QLoRA} will generate a largely beneficial effect by dramatically lowering barriers to finetuning advanced LLMs, making this technology far more broadly accessible.